\chapter{§4.4 频率分布直方图与频率分布表}
\label{sec:4.4}


\prerequisite{§4.1 数据的收集与整理, §4.2 平均数、中位数与众数}

\gradelevel{7-8 年级}


当数据量很大时，逐个列出所有数据既不方便也看不出规律。我们需要把数据\textbf{分组}，然后观察各组中数据出现的\textbf{频率}，从而了解数据的整体\textbf{分布}情况。本章学习频率分布表和频率分布直方图的编制与解读。



\section*{频率与频次}


\begin{explore}


某校对八年级全体 200 名学生进行了一次体能测试，测得引体向上的成绩如下（部分数据）：

\end{explore}

\begin{quote}
5, 8, 12, 3, 7, 10, 6, 9, 11, 4, 8, 7, 6, 10, 5, ...（共 200 个数据）
\end{quote}


面对 200 个数据，直接观察几乎看不出任何规律。这些数据是集中在某个范围，还是分散得很开？大部分同学能做几个？

\begin{understand}


\textbf{频次}（也叫频数）：某个数值或某个区间内的数据出现的\textbf{次数}。

\textbf{频率}：某个数值或某个区间内的数据出现的次数占\textbf{总数据个数}的比值。

\[
\text{频率} = \frac{\text{频次}}{\text{数据总个数}}
\]


\textbf{性质}：
\begin{itemize}
  \item 每个区间的频率 $\geq 0$ 。
  \item 所有区间的频率之和 $= 1$ （即 $100\%$ ）。
\end{itemize}


例如，如果 200 名同学中有 36 人的引体向上成绩在 $[6, 8)$ 这个区间内，则该区间的频次 $= 36$ ，频率 $= \frac{36}{200} = 0.18 = 18\%$ 。

\end{understand}

\begin{keytakeaway}


\begin{itemize}
  \item 频次是数据落入某区间的\textbf{次数}，频率是频次占总数的\textbf{比例}。
  \item 所有区间的频率之和等于 $1$ 。
\end{itemize}


\end{keytakeaway}



\section*{频率分布表的编制}


\begin{explore}


继续上面的体能测试数据。为了了解全年级引体向上成绩的分布，老师决定把数据分组来分析。应该怎么分组？每组多宽合适？

\end{explore}

\begin{understand}


编制频率分布表的步骤：

\textbf{第一步：确定数据的范围。}

找出最大值和最小值，计算极差。假设最小值 $= 1$ ，最大值 $= 15$ ，则极差 $= 15 - 1 = 14$ 。

\textbf{第二步：确定组距和组数。}

\begin{itemize}
  \item \textbf{组距}：每组的宽度。
  \item \textbf{组数}：分成几组。通常取 $5 \sim 8$ 组。
  \item 组距 $= \frac{\text{极差}}{\text{组数}}$ 。
\end{itemize}


这里取组距 $= 3$ ，则需要 $\lceil 14 \div 3 \rceil = 5$ 组。

\textbf{第三步：确定各组的区间。}

为了不重不漏，采用\textbf{左闭右开}区间（最后一组可以左闭右闭）：

\[
[1, 4), \quad [4, 7), \quad [7, 10), \quad [10, 13), \quad [13, 16]
\]


\textbf{第四步：逐个统计各区间的频次。} 用"正"字计数法或划记法。

\textbf{第五步：计算各区间的频率。}

假设统计结果如下：

\begin{center}
\begin{tabular}{lll}
\toprule
成绩区间 & 频次 & 频率 \\
\midrule
$[1, 4)$ & 20 & $\frac{20}{200} = 0.10$ \\
$[4, 7)$ & 52 & $\frac{52}{200} = 0.26$ \\
$[7, 10)$ & 68 & $\frac{68}{200} = 0.34$ \\
$[10, 13)$ & 44 & $\frac{44}{200} = 0.22$ \\
$[13, 16]$ & 16 & $\frac{16}{200} = 0.08$ \\
\textbf{合计} & \textbf{200} & \textbf{1.00} \\
\bottomrule
\end{tabular}
\end{center}


验证：频次之和 $= 200$ ✓，频率之和 $= 1.00$ ✓。

\textbf{分组的注意事项}：
\begin{itemize}
  \item 各区间\textbf{不重不漏}，每个数据恰好属于一个区间。
  \item 组距要适当：太大则看不出分布细节，太小则每组数据太少，看不出规律。
  \item 通常取 $5 \sim 8$ 组为宜。
\end{itemize}


\end{understand}

\begin{workedexamples}


\textbf{例 1.} 某班 40 名同学的 100 米跑成绩（单位：秒）如下：

\end{workedexamples}

\begin{quote}
14.2, 15.0, 13.8, 16.1, 14.5, 15.3, 13.5, 14.8, 15.6, 14.0,
13.9, 15.2, 14.7, 16.0, 13.6, 14.3, 15.1, 14.9, 15.5, 13.7,
14.1, 15.4, 14.6, 16.2, 13.4, 14.4, 15.8, 14.2, 15.0, 13.8,
14.5, 15.3, 13.6, 14.8, 15.6, 14.0, 13.9, 15.2, 14.7, 16.0
\end{quote}


编制频率分布表（组距取 0.5 秒）。

\textbf{解：}

第一步：最小值 $= 13.4$ ，最大值 $= 16.2$ ，极差 $= 16.2 - 13.4 = 2.8$ 。

第二步：组距 $= 0.5$ ，组数 $= \lceil 2.8 \div 0.5 \rceil = 6$ 组。

第三步：分组为 $[13.4, 13.9)$ , $[13.9, 14.4)$ , $[14.4, 14.9)$ , $[14.9, 15.4)$ , $[15.4, 15.9)$ , $[15.9, 16.4]$ 。

第四步和第五步：逐个统计。

\begin{center}
\begin{tabular}{lll}
\toprule
成绩区间（秒） & 频次 & 频率 \\
\midrule
$[13.4, 13.9)$ & 7 & $\frac{7}{40} = 0.175$ \\
$[13.9, 14.4)$ & 8 & $\frac{8}{40} = 0.200$ \\
$[14.4, 14.9)$ & 8 & $\frac{8}{40} = 0.200$ \\
$[14.9, 15.4)$ & 8 & $\frac{8}{40} = 0.200$ \\
$[15.4, 15.9)$ & 5 & $\frac{5}{40} = 0.125$ \\
$[15.9, 16.4]$ & 4 & $\frac{4}{40} = 0.100$ \\
\textbf{合计} & \textbf{40} & \textbf{1.000} \\
\bottomrule
\end{tabular}
\end{center}


验证： $7 + 8 + 8 + 8 + 5 + 4 = 40$ ✓

从频率分布表可以看出，大部分同学的成绩集中在 $[13.9, 15.4)$ 这三个区间，占全班的 $0.200 + 0.200 + 0.200 = 60\%$ 。

\begin{keytakeaway}


\begin{itemize}
  \item 频率分布表编制的五步：确定范围 $\to$ 定组距和组数 $\to$ 确定区间 $\to$ 统计频次 $\to$ 计算频率。
  \item 区间采用"左闭右开"原则，确保数据不重不漏。
  \item 频次之和 $=$ 数据总个数，频率之和 $= 1$ 。
\end{itemize}


\end{keytakeaway}



\section*{频率分布直方图的绘制与解读}


\begin{explore}


频率分布表给出了精确的数字，但当我们需要\textbf{一眼看出}数据分布的形态（集中在哪里？是对称还是偏斜？）时，一张图远比一张表直观。

\end{explore}

\begin{understand}


\textbf{频率分布直方图}是将频率分布表中的信息用图形表示出来的一种统计图。

\textbf{绘制方法}：

\begin{enumerate}
  \item 画出横轴（数据区间）和纵轴。
  \item 纵轴表示 $\frac{\text{频率}}{\text{组距}}$ （即\textbf{频率/组距}，也叫频率密度）。
  \item 在每个区间上画一个矩形，矩形的\textbf{宽度}等于组距，\textbf{高度}等于该区间的频率/组距。
  \item 因此，每个矩形的\textbf{面积} $=$ 组距 $\times \frac{\text{频率}}{\text{组距}} = \text{频率}$ 。
\end{enumerate}


\textbf{为什么纵轴是频率/组距而不是频率？}

如果各组的组距不同，直接用频率作为高度会产生误导——宽度大的矩形视觉上面积大，但频率不一定大。用频率/组距作为高度，保证了\textbf{矩形面积 $=$ 频率}，这样视觉上面积大的区间确实包含了更多的数据。

\textbf{当组距相同时}（大多数情况下），纵轴也可以直接标频率（因为频率与频率/组距成正比），但标准做法是标频率/组距。

\textbf{解读频率分布直方图的要点}：
\begin{itemize}
  \item \textbf{整体形态}：数据集中在哪个区间？是左右对称还是偏向一边？
  \item \textbf{峰值}：最高的矩形对应数据最集中的区间。
  \item \textbf{面积意义}：某几个相邻矩形的面积之和 $=$ 数据落在这些区间内的频率。
  \item \textbf{所有矩形面积之和} $= 1$ 。
\end{itemize}


\end{understand}

\begin{workedexamples}


\textbf{例 2.} 根据例 1 的频率分布表，绘制频率分布直方图。

\textbf{解：}

组距 $= 0.5$ ，计算各区间的频率/组距：

\begin{center}
\begin{tabular}{lll}
\toprule
成绩区间（秒） & 频率 & 频率/组距 \\
\midrule
$[13.4, 13.9)$ & $0.175$ & $\frac{0.175}{0.5} = 0.35$ \\
$[13.9, 14.4)$ & $0.200$ & $\frac{0.200}{0.5} = 0.40$ \\
$[14.4, 14.9)$ & $0.200$ & $\frac{0.200}{0.5} = 0.40$ \\
$[14.9, 15.4)$ & $0.200$ & $\frac{0.200}{0.5} = 0.40$ \\
$[15.4, 15.9)$ & $0.125$ & $\frac{0.125}{0.5} = 0.25$ \\
$[15.9, 16.4]$ & $0.100$ & $\frac{0.100}{0.5} = 0.20$ \\
\bottomrule
\end{tabular}
\end{center}


以成绩（秒）为横轴，频率/组距为纵轴，在每个区间上画出对应高度的矩形，各矩形紧密相邻（无间隔），即得频率分布直方图。

验证：所有矩形面积之和 $= 0.175 + 0.200 + 0.200 + 0.200 + 0.125 + 0.100 = 1.00$ ✓

\begin{figure}[htbp]
\centering
\begin{tikzpicture}
\begin{axis}[
  ybar interval,
  xmin=13.4, xmax=16.4,
  xtick={13.4,13.9,14.4,14.9,15.4,15.9,16.4},
  xticklabel style={rotate=45, anchor=east, font=\small},
  xlabel={成绩（秒）},
  ylabel={频率/组距},
  ymin=0, ymax=0.50,
  width=10cm, height=6cm,
  bar width=0.5,
  fill=blue!25, draw=blue!60,
]
\addplot coordinates {
  (13.4,0.35)(13.9,0.40)(14.4,0.40)(14.9,0.40)(15.4,0.25)(15.9,0.20)(16.4,0)
};
\end{axis}
\end{tikzpicture}
\caption{百米跑成绩频率分布直方图（组距 $= 0.5$ 秒）}
\label{fig:freq-hist-sprint}
\end{figure}

\textbf{例 3.} 下面的频率分布直方图表示某次考试 50 名学生的成绩分布（组距为 10 分）：

\begin{center}
\begin{tabular}{ll}
\toprule
成绩区间 & 频率/组距 \\
\midrule
$[40, 50)$ & $0.004$ \\
$[50, 60)$ & $0.012$ \\
$[60, 70)$ & $0.020$ \\
$[70, 80)$ & $0.036$ \\
$[80, 90)$ & $0.024$ \\
$[90, 100]$ & $0.004$ \\
\bottomrule
\end{tabular}
\end{center}

\begin{figure}[htbp]
\centering
\begin{tikzpicture}
\begin{axis}[
  ybar interval,
  xmin=40, xmax=100,
  xtick={40,50,60,70,80,90,100},
  xlabel={成绩（分）},
  ylabel={频率/组距},
  ymin=0, ymax=0.045,
  width=10cm, height=6cm,
  bar width=10,
  fill=green!25, draw=green!60,
]
\addplot coordinates {
  (40,0.004)(50,0.012)(60,0.020)(70,0.036)(80,0.024)(90,0.004)(100,0)
};
\end{axis}
\end{tikzpicture}
\caption{某次考试成绩频率分布直方图（组距 $= 10$ 分，$n=50$）}
\label{fig:freq-hist-exam}
\end{figure}

（1）成绩在 $[70, 80)$ 区间的频率是多少？有多少人？
（2）成绩在 $70$ 分以上（含 $70$ 分）的学生占多大比例？

\textbf{解：}

先验证数据：所有频率之和 $= (0.004 + 0.012 + 0.020 + 0.036 + 0.024 + 0.004) \times 10 = 0.100 \times 10 = 1.00$ ✓

（1）成绩在 $[70, 80)$ 区间的频率 $=$ 频率/组距 $\times$ 组距 $= 0.036 \times 10 = 0.36$ ，即 $36\%$ 。

人数 $= 50 \times 0.36 = 18$ 人。

（2）成绩在 $70$ 分以上的频率：

\[
0.036 \times 10 + 0.024 \times 10 + 0.004 \times 10 = 0.36 + 0.24 + 0.04 = 0.64
\]


占 $64\%$ 。

\end{workedexamples}

\begin{keytakeaway}


\begin{itemize}
  \item 频率分布直方图中，每个矩形的\textbf{面积 $=$ 频率}，所有面积之和 $= 1$ 。
  \item 纵轴标注\textbf{频率/组距}，保证面积具有频率意义。
  \item 从直方图可以直观看出数据集中的区间和分布形态。
\end{itemize}


\end{keytakeaway}



\section*{用频率分布直方图估计统计量}


\begin{explore}


如果我们只有一张频率分布直方图（原始数据已经丢失），还能估计出这组数据的平均数、中位数和众数吗？

\end{explore}

\begin{understand}


\textbf{估计众数}：

众数大约位于\textbf{最高矩形}所对应区间的\textbf{中点}。

在例 3 中，最高矩形对应 $[70, 80)$ ，该区间中点为 $\frac{70 + 80}{2} = 75$ ，因此估计众数约为 $75$ 分。

\textbf{估计中位数}：

中位数将所有数据分成\textbf{相等的两半}。在直方图中，中位数就是使得左边所有矩形面积之和 $= 0.5$ 的那条竖线的位置。

从例 3 的数据：
\begin{itemize}
  \item $[40, 50)$ 的频率 $= 0.04$
  \item $[50, 60)$ 的频率 $= 0.12$
  \item $[60, 70)$ 的频率 $= 0.20$
  \item 前三组频率之和 $= 0.04 + 0.12 + 0.20 = 0.36 < 0.5$
\end{itemize}


中位数在 $[70, 80)$ 区间内。还需 $0.5 - 0.36 = 0.14$ ，而 $[70, 80)$ 的频率 $= 0.36$ 。

\[
\text{中位数} \approx 70 + \frac{0.14}{0.36} \times 10 \approx 70 + 3.9 \approx 73.9 \text{（分）}
\]


\textbf{估计平均数}：

假设每个区间内的数据都集中在区间\textbf{中点}，用加权平均数估算。

\[
\begin{aligned}
\bar{x} &\approx 45 \times 0.04 + 55 \times 0.12 + 65 \times 0.20 + 75 \times 0.36 + 85 \times 0.24 + 95 \times 0.04 \\
&= 1.8 + 6.6 + 13.0 + 27.0 + 20.4 + 3.8 \\
&= 72.6 \text{（分）}
\end{aligned}
\]


\end{understand}

\begin{workedexamples}


\textbf{例 4.} 某次体检中，60 名学生的身高数据整理成如下频率分布表：

\begin{center}
\begin{tabular}{lll}
\toprule
身高区间（cm） & 频次 & 频率 \\
\midrule
$[145, 150)$ & 3 & $0.05$ \\
$[150, 155)$ & 9 & $0.15$ \\
$[155, 160)$ & 18 & $0.30$ \\
$[160, 165)$ & 21 & $0.35$ \\
$[165, 170)$ & 6 & $0.10$ \\
$[170, 175]$ & 3 & $0.05$ \\
\textbf{合计} & \textbf{60} & \textbf{1.00} \\
\bottomrule
\end{tabular}
\end{center}


（1）绘制频率分布直方图（组距为 5 cm）。
（2）估计这组数据的众数、中位数和平均数。

\textbf{解：}

（1）计算频率/组距：

\begin{center}
\begin{tabular}{ll}
\toprule
身高区间（cm） & 频率/组距 \\
\midrule
$[145, 150)$ & $\frac{0.05}{5} = 0.010$ \\
$[150, 155)$ & $\frac{0.15}{5} = 0.030$ \\
$[155, 160)$ & $\frac{0.30}{5} = 0.060$ \\
$[160, 165)$ & $\frac{0.35}{5} = 0.070$ \\
$[165, 170)$ & $\frac{0.10}{5} = 0.020$ \\
$[170, 175]$ & $\frac{0.05}{5} = 0.010$ \\
\bottomrule
\end{tabular}
\end{center}


以身高（cm）为横轴，频率/组距为纵轴，画出各矩形即可。

（2）

\textbf{众数}：最高矩形对应 $[160, 165)$ ，估计众数 $\approx \frac{160 + 165}{2} = 162.5$ cm。

\textbf{中位数}：
\begin{itemize}
  \item 前两组频率之和 $= 0.05 + 0.15 = 0.20$
  \item 前三组频率之和 $= 0.20 + 0.30 = 0.50 = 0.5$
\end{itemize}


恰好前三组的频率之和 $= 0.5$ ，所以中位数在第三组 $[155, 160)$ 的右端点，即中位数 $\approx 160$ cm。

\textbf{平均数}：

\[
\begin{aligned}
\bar{x} &\approx 147.5 \times 0.05 + 152.5 \times 0.15 + 157.5 \times 0.30 + 162.5 \times 0.35 + 167.5 \times 0.10 + 172.5 \times 0.05 \\
&= 7.375 + 22.875 + 47.25 + 56.875 + 16.75 + 8.625 \\
&= 159.75 \text{（cm）}
\end{aligned}
\]


\end{workedexamples}

\begin{keytakeaway}


\begin{itemize}
  \item 从频率分布直方图估计\textbf{众数}：取最高矩形的区间中点。
  \item 从频率分布直方图估计\textbf{中位数}：找使左侧面积 $= 0.5$ 的位置。
  \item 从频率分布直方图估计\textbf{平均数}：以各区间中点为代表值，用频率作为权重计算加权平均数。
  \item 这些都是\textbf{估计值}，不一定等于真实的统计量。
\end{itemize}


\end{keytakeaway}

\begin{practice}


\textbf{1.} 某校 100 名八年级学生 1 分钟跳绳成绩整理如下：

\begin{center}
\begin{tabular}{ll}
\toprule
成绩区间（次） & 频次 \\
\midrule
$[60, 80)$ & 8 \\
$[80, 100)$ & 22 \\
$[100, 120)$ & 35 \\
$[120, 140)$ & 25 \\
$[140, 160]$ & 10 \\
\bottomrule
\end{tabular}
\end{center}


（1）补全频率分布表（计算频率）。
（2）计算各区间的频率/组距，并描述如何绘制频率分布直方图。
（3）成绩在 $100$ 次以上（含 $100$ 次）的学生占百分之几？

\textbf{2.} 根据下面的频率分布表，估计这组数据的众数、中位数和平均数。

\begin{center}
\begin{tabular}{ll}
\toprule
数据区间 & 频率 \\
\midrule
$[10, 20)$ & $0.10$ \\
$[20, 30)$ & $0.20$ \\
$[30, 40)$ & $0.40$ \\
$[40, 50)$ & $0.20$ \\
$[50, 60]$ & $0.10$ \\
\bottomrule
\end{tabular}
\end{center}


\textbf{3.} 某频率分布直方图中， $[20, 25)$ 区间的频率/组距为 $0.08$ ，组距为 $5$ 。
（1）该区间的频率是多少？
（2）如果数据共有 200 个，该区间有多少个数据？

\end{practice}



\begin{answer}


\textbf{1.}

（1）总人数 $= 100$ 。

\begin{center}
\begin{tabular}{lll}
\toprule
成绩区间（次） & 频次 & 频率 \\
\midrule
$[60, 80)$ & 8 & $0.08$ \\
$[80, 100)$ & 22 & $0.22$ \\
$[100, 120)$ & 35 & $0.35$ \\
$[120, 140)$ & 25 & $0.25$ \\
$[140, 160]$ & 10 & $0.10$ \\
\textbf{合计} & \textbf{100} & \textbf{1.00} \\
\bottomrule
\end{tabular}
\end{center}


（2）组距 $= 20$ 。频率/组距依次为：

\[
\frac{0.08}{20} = 0.004, \quad \frac{0.22}{20} = 0.011, \quad \frac{0.35}{20} = 0.0175, \quad \frac{0.25}{20} = 0.0125, \quad \frac{0.10}{20} = 0.005
\]


以成绩为横轴，频率/组距为纵轴，在每个区间上画出对应高度的矩形，各矩形紧密相邻。

（3）成绩 $\geq 100$ 的频率 $= 0.35 + 0.25 + 0.10 = 0.70$ ，即 $70\%$ 。

\textbf{2.}

\textbf{众数}：最高频率的区间是 $[30, 40)$ （频率 $= 0.40$ ），众数 $\approx \frac{30+40}{2} = 35$ 。

\textbf{中位数}：
\begin{itemize}
  \item $[10, 20)$ 频率 $= 0.10$
  \item $[10, 20) + [20, 30)$ 频率 $= 0.10 + 0.20 = 0.30 < 0.5$
\end{itemize}


中位数在 $[30, 40)$ 区间。还需 $0.5 - 0.30 = 0.20$ ， $[30, 40)$ 的频率 $= 0.40$ 。

\[
\text{中位数} \approx 30 + \frac{0.20}{0.40} \times 10 = 30 + 5 = 35
\]


\textbf{平均数}：

\[
\bar{x} \approx 15 \times 0.10 + 25 \times 0.20 + 35 \times 0.40 + 45 \times 0.20 + 55 \times 0.10
\]


\[
= 1.5 + 5.0 + 14.0 + 9.0 + 5.5 = 35
\]


\textbf{3.}

（1）频率 $= 0.08 \times 5 = 0.40$ ，即 $40\%$ 。

（2）数据个数 $= 200 \times 0.40 = 80$ 个。

\end{answer}
